% 第 11 题程序框图（含「输入」框）：按固定源 PDF 矢量路径逐框重绘（1 pt → 0.03515 cm）。
% 原点取主干竖线（源页 x 564.62）与「结束」底边（y 680.96），y 向上。共 9 框 9 箭头，全实线。
% 结构要点（易错）：「输入 n」与「输出 S」同为向右倾平行四边形；处理框「S=S+√(i+1)−√i」为全图最宽；
% i=i+1 在右侧独立框；「否」支自菱形右顶点水平右行、折向上入 i=i+1 底边；i=i+1 的出口在其【上边】，
% 向上至 y=121.28 后折左汇入主干——汇入点在 S=0 与处理框之间，故循环体含处理框整段（非尾循环）。
% 初值 i=1、S=0；判断框内为单线大于等于号（U+2A7E 斜线双划），照源。
\begin{tikzpicture}[>=stealth,every node/.style={font=\small},x=0.03515cm,y=0.03515cm]
% ===== 框 =====
\draw[rounded corners=3.98pt] (-13.99,217.18) rectangle (13.99,234.02);        % 开始
\draw (-25.02,185.78)--(18.89,185.78)--(25.01,202.62)--(-18.90,202.62)--cycle; % 输入 n（右倾平行四边形）
\draw (-14.93,157.30) rectangle (14.93,171.21);    % i=1
\draw (-16.64,128.56) rectangle (16.64,142.73);    % S=0
\draw (-51.77,95.64) rectangle (51.77,113.99);     % S=S+√(i+1)−√i（最宽）
\draw (66.33,97.41) rectangle (112.63,112.21);      % i=i+1（右侧独立）
\draw (36.05,72.04)--(0,80.96)--(-36.06,72.04)--(0,63.09)--cycle; % 菱形 i≥n
\draw (-25.40,31.59)--(19.26,31.59)--(25.39,48.44)--(-19.27,48.44)--cycle; % 输出 S（右倾平行四边形）
\draw[rounded corners=3.98pt] (-13.99,0) rectangle (13.99,17.02);              % 结束
% ===== 框内文字 =====
\node at (0,226.11) {开始};
\node at (0,194.5) {输入 \(n\)};
\node at (0,163.33) {\(i=1\)};
\node at (0,134.58) {\(S=0\)};
\node at (0,108.31) {\(S=S+\sqrt{i+1}-\sqrt{i}\)};
\node at (89.48,104.31) {\(i=i+1\)};
\node at (0,72.33) {\(i\geqslant n\)};
\node at (0,40.00) {输出 \(S\)};
\node at (0,8.97) {结束};
% ===== 箭头（主干七段 + 否支 + 回接）=====
\draw[->] (0,233.82)--(0,202.62);   % 开始→输入
\draw[->] (0,185.78)--(0,171.21);   % 输入→i=1
\draw[->] (0,157.30)--(0,142.73);   % i=1→S=0
\draw[->] (0,128.56)--(0,113.99);   % S=0→处理框
\draw[->] (0,95.64)--(0,80.96);     % 处理框→菱形
\draw[->] (0,63.09)--(0,48.44);     % 菱形→输出（是分支）
\draw[->] (0,31.59)--(0,17.02);     % 输出→结束
\draw[->] (36.05,72.04)--(89.48,72.04)--(89.48,97.41);  % 否支：右顶点→右→上入 i=i+1 底边
\draw[->] (89.48,112.21)--(89.48,121.28)--(0,121.28);   % 回接：i=i+1 上边→上→左汇入主干
% ===== 分支标注 =====
\node at (9.16,56.28) {是};
\node at (36.56,64.20) {否};
\end{tikzpicture}
